% The Elsevier (elsarticle) body of the note, shared by the two wrappers:
%   note_ipl.tex        -- [review,12pt]: double spaced with line numbers, for submission
%   note_ipl_print.tex  -- [final,3p,times,twocolumn]: the journal's print layout, used only
%                          to check the length against the nine-printed-page limit
% The mathematics and the text are not duplicated: both wrappers include the same
% ../../note/note-abstract.tex and ../../note/note-content.tex as the standalone note does,
% so scripts/w_checknums.py still holds every number in one place.
\usepackage{amsmath,amssymb,amsthm}
\usepackage{booktabs}
\usepackage[hidelinks]{hyperref}

\newtheorem{theorem}{Theorem}
\newtheorem{lemma}[theorem]{Lemma}
\newtheorem{proposition}[theorem]{Proposition}
\newtheorem{corollary}[theorem]{Corollary}
\theoremstyle{definition}
\newtheorem{definition}[theorem]{Definition}
\theoremstyle{remark}
\newtheorem{observation}[theorem]{Observation}
\newtheorem{remark}[theorem]{Remark}

\newcommand{\PH}{P^{\mathrm H}}
\newcommand{\PV}{P^{\mathrm V}}
\newcommand{\Zs}{\mathbb{Z}_7^5}
\newcommand{\boxt}{\boxtimes}

\journal{Information Processing Letters}

\begin{document}

\begin{frontmatter}

\title{The Polak--Schrijver code has exactly eight private pairs}

\author{Artem Oktiabrev}
\ead{aoktyabrev@gmail.com}
\ead[url]{https://orcid.org/0009-0003-3626-2002}

\begin{abstract}
% The abstract, shared by note.tex and submission/ipl/note_ipl.tex.
Every improvement on the lower bound for the Shannon capacity $\Theta(C_7)$ since July 2026
is built on one five-dimensional gadget: the $367$-word independent set of Polak and
Schrijver together with eight \emph{private pairs} --- vertices outside the code with a
single neighbour in it --- and the private-pair count is the parameter those recursions are
most sensitive to. We prove that eight is the exact maximum: in the whole of $\Zs$ exactly
eight vertices have a single neighbour in the code, they are precisely the eight in use, and
all eight are simultaneously usable because their confusability graph is a matching, hence
bipartite. A ninth private pair for this code does not exist. We also report the exhaustive
computations around the theorem: all eight codes the Polak--Schrijver pipeline can produce,
with counts $5,6,6,6,7,7,7,8$; the arithmetic that rules out smaller codes; and a sweep of
$8{,}260{,}976{,}640$ images under $\mathrm{Aut}(C_7^{\boxt 5})$ which shows why the
published construction has to replace one auxiliary vector.

\end{abstract}

\begin{keyword}
Shannon capacity \sep zero-error capacity \sep independent set \sep strong product \sep
odd cycle \sep gadget product
\MSC[2020] 05C69 \sep 05C76 \sep 94A24
\end{keyword}

\end{frontmatter}

% Sections 1-7 and the closing paragraphs, shared by note.tex (standalone,
% archived with the Zenodo record) and submission/ipl/note_ipl.tex (elsarticle,
% for submission).  Every number in this file is checked by scripts/w_checknums.py.
\section{Context, and what the question is worth}
\label{sec:intro}

For a graph $G$, the Shannon capacity is $\Theta(G)=\sup_d \alpha(G^{\boxt d})^{1/d}$
\cite{shannon1956}. For the seven-cycle it is unknown, and the gap between the best bounds
is wide: the Lov\'asz bound gives
$\Theta(C_7)\le\vartheta(C_7)=3.3176672073940953927\ldots$ \cite{lovasz1979}, while the
lower bound has moved only in the fourth decimal place in seven years. Table~\ref{tab:chain}
lists the chain.

\begin{table}[ht]
\centering\small
\begin{tabular}{@{}llll@{}}
\toprule
date & $\Theta(C_7)\ge$ & from & witness \\
\midrule
1971--2015 & $343^{1/5}$, then $350^{1/5}$ & \cite{baumert1971,mathew2017} & a linear code; stochastic search \\
2019 & $3.2578659667835159504124\ldots$ & \cite{polak2019} & $\alpha(C_7^{\boxt 5})\ge 367$ \\
Jul 2026 & $3.2580207372932453595278\ldots$ & \cite{itty2026} & $\alpha(C_7^{\boxt 10})\ge 134753$ \\
Jul 2026 & $3.2587891539086910161967650155206769\ldots$ & \cite{gao2026} & recursive gadget product \\
Jul 2026 & $3.2588053698854655725829750306238067\ldots$ & \cite{bpz2026} & a stronger base gadget \\
Aug 2026 & $3.2588326203532663091215390518104754\ldots$ & \cite{tandon2026} & heterogeneous recursion \\
\bottomrule
\end{tabular}
\caption{Lower bounds on $\Theta(C_7)$. The last four entries all start from the same
$367$-word code and the same eight private pairs; they differ in the recursion built on
top.}
\label{tab:chain}
\end{table}

The structure behind the last four rows is Gao's gadget \cite{gao2026}, restated as a
\emph{valid tuple} in \cite{bpz2026} and refined in \cite{tandon2026}: an independent set
$I$ (the code), $t$ private pairs, two complementary transversals of their endpoints, and an
auxiliary independent set $X$. The product lemma multiplies gadgets, and the parameter
sensitivity of the resulting bound is lopsided. Evaluated at hypothetical base profiles
$(a,t,s,o)$, the homogeneous recursion of \cite{gao2026,bpz2026} at its optimal depth gives
\[
(367,8,367,322)\mapsto 3.2588053698854655725829750306238067\ldots,\qquad
(367,9,367,321)\mapsto 3.2590062007738521305802\ldots,
\]
so one further private pair is worth $1.736\cdot 10^{-4}$ over the current record, whereas one
further auxiliary word of the neutral class is worth about a tenth of that
($(367,8,367,324)$, two extra such words, clears the record by only
$6.993\cdot10^{-6}$). For scale: the difference between the record \cite{tandon2026} and what
the homogeneous recursion extracts from the \emph{same} base gadget is $2.725\cdot10^{-5}$, so
a ninth private pair would improve the state of the art under any of the published recursions,
by roughly seven times the spread between the published schemes. That is why the question
of a ninth pair is worth settling rather than leaving as folklore, and it is the question this
note answers.

The answer is negative and exact, which also explains a pattern in the literature: the three
independent lines of work after \cite{itty2026} all keep the same base gadget and improve the
recursion around it. On that code the private-pair count cannot be raised at all.

\section{Gadgets and private pairs}
\label{sec:defs}

Vertices of $C_7^{\boxt d}$ are the elements of $\mathbb{Z}_7^d$; two distinct vertices are
adjacent (\emph{confusable}) iff their circular distance is at most $1$ in every coordinate.
We write $N[v]$ for the closed neighbourhood and call $I\subseteq\mathbb{Z}_7^d$ independent
if no two of its elements are confusable.

\begin{definition}[private pair, \cite{gao2026}]
\label{def:pair}
Let $I$ be independent. A pair $(r,q)$ with $r\in I$, $q\notin I$ is a \emph{private pair}
for $I$ if $N[q]\cap I=\{r\}$. The vertex $r$ is the \emph{centre} and $q$ the \emph{private
neighbour}.
\end{definition}

\begin{definition}[gadget, \cite{gao2026}]
\label{def:gadget}
A gadget consists of an independent set $I$ of size $a$; private pairs $(r_i,q_i)$,
$1\le i\le t$, with pairwise disjoint endpoint sets; two complementary transversals
$\PH,\PV$ of the endpoints --- for every $i$ exactly one of $r_i,q_i$ lies in $\PH$ and the
other in $\PV$ --- both independent; and an independent auxiliary set $X$ with
$X\cap N(\PH)\cap N(\PV)=\emptyset$. With
$O=X\setminus(N(\PH)\cup N(\PV))$, $H=X\cap N(\PH)$, $V=X\cap N(\PV)$, the profile is
$(a,t,s,o,h,v)=(a,t,|X|,|O|,|H|,|V|)$.
\end{definition}

Write $t^*(I)$ for the largest $t$ for which some gadget with code $I$ exists. The published
base gadget has $I=I_0$, the $367$-word set printed in the appendix of \cite{polak2019},
and $t=8$ with the eight pairs of \cite{itty2026}; its profile is $(367,8,367,321,26,20)$
in \cite{gao2026} and $(367,8,367,322,26,19)$ in \cite{bpz2026,tandon2026}.

\section{The theorem}
\label{sec:theorem}

The definition of a gadget quantifies over transversals, which is what makes $t^*$ look hard
to bound. It is not: the transversal condition collapses to a condition on the private
neighbours alone.

\begin{lemma}[no mixing]
\label{lem:nomix}
Let $(r_i,q_i)$ and $(r_j,q_j)$ be private pairs for $I$ with $r_i\neq r_j$. Then $r_i$ is
confusable with neither $r_j$ nor $q_j$.
\end{lemma}

\begin{proof}
$r_i,r_j\in I$ and $I$ is independent, so $r_i\not\sim r_j$. If $r_i\sim q_j$ then
$r_i\in N[q_j]\cap I=\{r_j\}$, so $r_i=r_j$, contradicting the hypothesis.
\end{proof}

\begin{corollary}
\label{cor:bip}
Let $(r_i,q_i)$, $1\le i\le t$, be private pairs for $I$ with pairwise distinct centres, and
let $\sigma:\{1,\dots,t\}\to\{\mathrm H,\mathrm V\}$ assign to each pair the transversal
holding $q_i$. Then the endpoint sets are pairwise disjoint, and $\PH,\PV$ are both
independent if and only if $\sigma$ is a proper $2$-colouring of the graph $\Gamma$ on
$\{q_1,\dots,q_t\}$ whose edges are the confusable pairs $q_i\sim q_j$. Consequently
$t^*(I)$ is the largest number of private pairs with distinct centres whose $q$-graph
$\Gamma$ is bipartite.
\end{corollary}

\begin{proof}
Distinct centres give $r_i\neq r_j$; $q_i\neq q_j$ because $N[q_i]\cap I=\{r_i\}$ pins the
centre; and $r_i\neq q_j$ because $r_i\in I$, $q_j\notin I$. So the endpoint sets are
disjoint. Each transversal consists of some $r$'s and some $q$'s. Two $r$'s are never
confusable, and an $r$ and a $q$ of different pairs are never confusable by
Lemma~\ref{lem:nomix}; an $r_i$ and $q_i$ of the \emph{same} pair are confusable but lie in
different transversals by construction. Hence a transversal can fail to be independent only
by containing two confusable private neighbours, which is exactly a monochromatic edge of
$\Gamma$.
\end{proof}

So the only candidates that need to be examined are the private pairs of $I$ themselves, and
they form a finite list which can be written down: every $q\notin I$ with $|N[q]\cap I|=1$,
together with that single neighbour. For the $367$-word code this list is short.

\begin{proposition}[the enumeration]
\label{prop:count}
Let $I_0\subseteq\Zs$ be the $367$-word Polak--Schrijver code as printed in \cite{polak2019}.
Of the $16807-367=16440$ vertices outside $I_0$, exactly $8$ have one neighbour in $I_0$,
$254$ have two, and $16178$ have three or more; none has zero. The eight are exactly the
private neighbours $q_0,\dots,q_7$ of \cite{itty2026,gao2026}, with the centres those papers
give them.
\end{proposition}

\begin{proof}
One pass over $\Zs$, counting for each vertex outside $I_0$ its neighbours in $I_0$; the
predicate is the integer condition of Section~\ref{sec:defs}, so the computation is exact.
The empty zero-bucket says that $I_0$ is a maximal independent set, which is independently
known \cite[Appendix]{polak2019}. The identification with the published pairs is a set
comparison. See Section~\ref{sec:repro}.
\end{proof}

\begin{theorem}
\label{thm:main}
$t^*(I_0)=8$. Moreover the admissible transversal assignments are exactly the proper
$2$-colourings of a graph with five connected components, i.e.\ $16$ of them up to exchanging
$\PH$ and $\PV$, and one of these is the assignment used in \cite{itty2026,gao2026,bpz2026,tandon2026}.
\end{theorem}

\begin{proof}
By Corollary~\ref{cor:bip}, $t^*(I_0)$ is computed from the candidate list of
Proposition~\ref{prop:count}. That list has eight entries with eight distinct centres, so
$t^*(I_0)\le 8$. Its conflict graph $\Gamma$ has, in the numbering of \cite{gao2026},
exactly the three edges
\[
  q_0\sim q_1,\qquad q_2\sim q_6,\qquad q_3\sim q_5,
\]
i.e.\ $\Gamma$ is a perfect matching on six of the eight vertices and has two isolated
vertices. A matching is bipartite, so all eight pairs may be used at once and
$t^*(I_0)=8$. The number of components of $\Gamma$ is $3+2=5$, and a graph whose components
are single edges and isolated vertices has $2^5$ proper $2$-colourings, i.e.\ $2^4=16$ up to
a global swap; that the published assignment ($J_0=\{0,5,6\}$ carrying $r_j$ into $\PH$,
$J_1=\{1,2,3,4,7\}$ the other way, as in \cite{gao2026}) is one of them is checked directly.
\end{proof}

\begin{corollary}
No gadget whose code is $I_0$ has $t\ge 9$. In particular the profile $(367,9,367,\cdot)$,
which would give $\Theta(C_7)\ge 3.2590062007738521305802\ldots$ through the homogeneous
recursion of \cite{gao2026,bpz2026}, is not reachable from the Polak--Schrijver code, and the
same holds for every recursion that takes a gadget over this code as its base, including
\cite{tandon2026}.
\end{corollary}

Lemma~\ref{lem:nomix} and Corollary~\ref{cor:bip} are proved by hand; the two finite facts
they are applied to --- the candidate list and the three edges --- are direct computations over
$16807$ vertices, small enough for any reader to repeat from the printed set.

\section{The exhaustive computations around the theorem}
\label{sec:exhaustive}

Theorem~\ref{thm:main} is about one code. Three further questions are settled by computation
rather than by proof, and together they close the directions in which a ninth pair could have
been sought.

\paragraph{All codes the Polak--Schrijver construction can produce.}
Step (v) of \cite{polak2019} extends a core of $327$ words by a maximum independent set of an
extension graph with $71$ vertices and $85$ edges, whose independence number is $40$. That
maximum is not unique. Enumerating all of them gives exactly $8$ distinct codes of size
$367$; their private-pair counts, each computed as in Theorem~\ref{thm:main}, are
\[
  5,\;6,\;6,\;6,\;7,\;7,\;7,\;8 ,
\]
and the unique code attaining $8$ is the printed one. The construction that produced the
record therefore has nothing better to give: the literature is already using the best member
of the family.

\paragraph{Codes of other sizes.}
A gadget's code need not be maximum, so a smaller code with more private pairs is admissible
in principle. The published recursion prices that trade exactly. Matching the bound of
$(367,8,367,321)$ would require
\[
  13 \text{ private pairs at } a=366,\qquad
  16 \text{ at } a=365,\qquad
  23 \text{ at } a=360 ,
\]
while the $367$-word code has eight candidates in the whole of $\Zs$. One lost code word
costs about six times what an extra private pair gains, so this branch is closed by
arithmetic and not by search.

\paragraph{The auxiliary set, over the whole automorphism group.}
With $t=8$ forced, the remaining parameter is $X$. Since $|X|$ should stay at $367$, the
candidates are the eight pipeline codes and their images under
$\mathrm{Aut}(C_7^{\boxt 5})\supseteq(D_7)^5\rtimes S_5$, a group of
$14^5\cdot 120=64{,}538{,}880$ elements; every admissible $2$-colouring changes $N(\PH)$ and
$N(\PV)$, so the sweep covers $8\cdot 16\cdot 64{,}538{,}880=8{,}260{,}976{,}640$ triples
(source code, colouring, automorphism). It is exhaustive, not sampled: for fixed permutation
and reflections all $16807$ translations are scored at once by a cyclic cross-correlation.
The result is that the forbidden intersection $X\cap N(\PH)\cap N(\PV)$ is never empty ---
the minimum over all $8{,}260{,}976{,}640$ triples is one word:

\begin{center}
\begin{tabular}{@{}lrrrrrrrr@{}}
\toprule
words of $X$ in the forbidden region & 0 & 1 & 2 & 3 & 4 & 5 & 6 & 7 \\
\midrule
triples & $0$ & $4$ & $33$ & $101$ & $209$ & $358$ & $965$ & $6{,}060$ \\
\bottomrule
\end{tabular}
\end{center}

This explains a step of the published construction that is stated there without comment: Gao
takes $X$ to be the image $T(I_0)$ under $T(w)=(2-w_1,w_3,w_0,2-w_2,w_4)$ with the single word
$(2,4,6,3,5)$ replaced by $(1,5,6,3,5)$ \cite{gao2026,tandon2026}. The replacement is
unavoidable rather than cosmetic, and $(2,4,6,3,5)$ is exactly the one word of $T(I_0)$
confusable with both transversals. Repairing the near misses exactly --- delete the offending
words, then solve for the best legal replacement --- reaches at best $s=367$, $o=322$: the
Buys--Polak--Zuiddam profile \cite{bpz2026}, recovered from another direction and not improved.
Across the $60$ repaired candidates no repair admitted more insertions than deletions; had one,
it would have been an independent set of size $368$ in $C_7^{\boxt 5}$, which the script watched
for and did not see.

\section{Two observations, with their status}
\label{sec:obs}

\begin{observation}[the colouring the literature uses may be the only usable one]
\label{obs:colour}
Of the $16$ proper $2$-colourings of Theorem~\ref{thm:main}, a $367$-word admissible
auxiliary set was found for exactly one --- the assignment of \cite{itty2026}, kept by
\cite{gao2026,bpz2026,tandon2026}. None of those papers remarks on the choice.
\end{observation}

The status of this observation is weaker than the theorem's, and is stated as such. The group
sweep is exhaustive and proves that no \emph{image} of a known code is admissible under any of the
$16$ colourings, so admissible sets have to be built by repair --- and repair is a search:
fixed-cardinality local search under each colouring for $1500$\,s reached $o=322$ under the
published colouring and nothing admissible under the other fifteen. For those fifteen the honest
statement is that $1500$\,s of search found nothing, not that nothing exists. The control that
makes the negative reading meaningful: started from the published auxiliary set under the
published colouring, the same search returns $o=321$, exactly the value of \cite{gao2026}.

\begin{observation}[the count is a property of the code, not of its size]
\label{obs:fragile}
An independent rebuild of the Polak--Schrijver construction, agreeing with the printed set in
$365$ of its $367$ words, admits only $6$ candidate private pairs; the linear $343$-word code
admits none at all.
\end{observation}

Two words out of $367$ change the count by two: nothing about the size, or about the
circular-graph construction, forces eight. That is the reason to look for a different count
outside this construction, and not to expect one within it.

\section{Reproducibility}
\label{sec:repro}

The computations above are exact integer computations, and the artefacts are published with
the note.

\emph{The verifier.} Independence is checked by a standalone C program with no dependencies.
For $n\ge4$ two distinct vertices of $C_n^{\boxt d}$ are adjacent exactly when the boxes
$\prod_i\{x_i,x_i+1\}$ meet, so a set is independent iff its $|I|\cdot 2^d$ box cells are all
distinct; the verifier sweeps them into an occupancy bitmap in one linear pass, in integer
arithmetic. It reports the size, the dimension, the verdict and the SHA-256 of the input file,
and it knows nothing about how a set was produced.

\emph{What the verifier is held to.} A verifier that always answered
\textsc{independent} would pass every reproduction, so the calibration is made of cases where
it can fail: $1177$ single-vertex corruptions against independently computed ground truth
(including $18$ that leave an independent set, which it must accept); maximality by exhaustive
addition over all $7^d$ vertices; the box method against a plain quadratic pair test at
$d\le4$; the chunked low-memory path against the single-pass path; and negative controls. A
last test asks whether those tests have teeth: a deliberately defective build of the same
verifier, seeing only repeated vertices, calls a dependent set independent and is rejected.

\emph{The sets.} Every set is published as one vertex per line with its SHA-256, in
\texttt{sets/SHA256SUMS}: the printed code (\texttt{9bdb8b9d69d9451e}), our rebuild of it
(\texttt{23826340a496cb04}), the eight pipeline codes, the best repaired auxiliary set
(\texttt{a4009726c06beb61}) and the linear $343$-word code (\texttt{fc4dcf8e7e214c8f}).

\emph{Seconds against hours.} Theorem~\ref{thm:main} and Proposition~\ref{prop:count} are a
matter of seconds: verifying a $367$-word set sweeps $11{,}744$ box cells, and the enumeration
is one pass over $16807$ vertices. (As a cross-check, the $134753$-word set of \cite{itty2026}
in $C_7^{\boxt 10}$ verifies in $0.47$\,s, $137{,}987{,}072$ cells.) What needs a full run is
the group sweep ($663.9$\,s) and the sixteen auxiliary searches ($1500$\,s each); the whole
study used about $18.7$ core-hours and no GPU time.

\emph{Independence from the code.} Theorem~\ref{thm:main} does not rest on the software: its
proof is Lemma~\ref{lem:nomix} plus two finite checks, each verifiable from the printed set in
\cite{polak2019} in a few lines of any language. The programs reproduce those lists; they are
not a premise of the argument.

\section{What remains open}
\label{sec:open}

\begin{enumerate}
\item \emph{A $367$-word code from outside the Polak--Schrijver construction.} The count is a
property of the individual code (Observation~\ref{obs:fragile}), and within the construction the
counts run $5,\dots,8$. A code with nine candidates would move the record by $1.7\cdot10^{-4}$;
producing one needs a new construction, or a direct attack on $\alpha(C_7^{\boxt 5})$, where
even $368$ is open.
\item \emph{A base gadget in a dimension other than five.} Every published gadget lives in
$C_7^{\boxt 5}$ because that is where the $367$ lives; nothing in the product lemma requires it.
\item \emph{The recursion, not the base.} The heterogeneous framework of \cite{tandon2026}
extracts $2.725\cdot10^{-5}$ more than the homogeneous one from the same base gadget, which is
more than the entire $o$-route on this code could give. Any future base improvement is worth
more through it.
\end{enumerate}

\paragraph{Data and code availability.}
The sets, the scripts that produce every number in this note, the machine reconciliation of
those numbers against the stored results, and the literature check are archived at
\href{https://doi.org/10.5281/zenodo.22979509}{doi:10.5281/zenodo.22979509}~\cite{shannonarchive}.

\paragraph{Declaration of competing interest.}
The author declares that he has no known competing financial interests or personal relationships
that could have appeared to influence the work reported in this paper.

\paragraph{Declaration of generative AI and AI-assisted technologies in the manuscript
preparation process.}
During the preparation of this work the author used Claude and Claude Code (Anthropic) in order
to draft and revise the text of this note. After using this tool/service the author reviewed
and edited the content as needed and takes full responsibility for the content of the published
article.

\paragraph{AI in the research process.}
The computations --- the enumeration of Proposition~\ref{prop:count}, the verifier and its
calibration, the pipeline family, the group sweep and the auxiliary searches --- were implemented
and run through Claude Code (Claude Opus 5) under the author's direction; Claude is a tool and not
an author. The safeguards are part of the process: predictions registered before each run, every
estimator calibrated on a case where it could fail (including a deliberately defective verifier
that had to be rejected), and every number in this note reconciled by script against the
computation behind it, an unsourced number counting as a failed build.


\bibliographystyle{elsarticle-num}
\bibliography{refs}

\end{document}
